\section{Hodge atoms and unconditional one stabilization}
\label{sec:atomic-one-step}

The proof of Theorem~\ref{thm:every-cubic} uses the ordinary Hodge-atom
birational ledger and one new finite invariant on a rank-two atom.  The
framed formal-monodromy theory is postponed to
Section~\ref{sec:framed-monodromy}; no reconstruction-displacement
hypothesis enters this section.

\subsection{Ordinary Hodge atoms and conventions}
We use the numbered ordinary Hodge-atom theorems of \cite{KKPYY} as external inputs.  Specifically we use the spectral decomposition of the
non-archimedean $A$-model $F$-bundle; Definition~5.21 and
Proposition~5.22, which associate a well-defined geometric atomic
$F$-bundle equivalence class to a Hodge atom; Proposition~5.23, which makes
the underlying $\Hod$-representation an atom invariant; the blowup and
projective-bundle chemical formulas used in the proof of Proposition~5.22;
Lemma~5.24 for varieties with nef canonical class; and the ordinary
non-rationality criterion, Proposition~5.30.  The projective-bundle part of
that package ultimately rests on the Iritani--Koto decomposition theorem, as
explained in \cite{KKPYY}.  At the time of writing, \cite{KKPYY} is an unrefereed preprint; the main theorem inherits that source risk.  No
enhanced atom, Euler pairing, Serre automorphism, or non-archimedean
Riemann--Hilbert assertion from \cite{KKPYY} is used here.

We also use Cai's small even quantum connection of a smooth cubic threefold
and his curve calculation \cite{Cai}.  The rank-two block data needed below
are reproduced explicitly, so the only imported cubic datum is the small
even connection matrix itself.

We use one standard rigid-analytic continuation fact.  In the form needed
below, if $Y$ is a connected normal rigid analytic space and an analytic
subset contains a nonempty admissible open, then it is all of $Y$; this is
\cite[Lemma~2.1.4]{Conrad}.  Smooth rigid spaces are normal.

\begin{lemma}[The intrinsic Euler endomorphism and the A-model spectral
operator]\label{lem:euler-sign}
\coverage{fragment}%
\lean{eulerOperator_generalizedEigenspaces_under_sign_reversal,%
  eulerOperator_eigenvalue_count_under_sign_reversal,%
  eulerOperator_cyclic_span_under_sign_reversal,%
  characteristicDiscriminant_under_root_sign_reversal,%
  residueDiscriminant_under_sign_reversal,%
  characteristicDiscriminant_under_odd_coefficient_sign_reversal}%
\imports{KKPYY:maximal-f-bundle, KKPYY:amodel-connection}%
Let $(\mathcal H,\nabla)/B$ be the maximal $A$-model $F$-bundle in the
conventions displayed in \cite{KKPYY}.  Let $E$ be the intrinsic Euler field
of the associated maximal $F$-manifold and let
\[
  \kappa:=\nabla_{u^2\partial_u}|_{u=0}
\]
be its residual endomorphism.  Then
\[
  C_E=\kappa,
  \qquad
  \kappa=-\Eu\star
\]
for the displayed $A$-model connection
$\nabla_{\partial_u}=\partial_u-u^{-2}(\Eu\star)+O(u^{-1})$.
Consequently $\kappa$ and $\Eu\star$ have the same generalized
eigenspace decomposition and the same number of distinct eigenvalues; their
reduced spectral covers are identified by $\lambda\mapsto-\lambda$.
Regularity of the endomorphism and the vanishing locus of its characteristic
polynomial discriminant are likewise unchanged.
\end{lemma}

\begin{proof}
By the maximal-$F$-bundle construction of \cite{KKPYY}, the intrinsic Euler
field is characterized by $C_E=\kappa$.  Reading the coefficient of
$u^{-2}$ in their displayed $A$-model connection gives
$\kappa=-\Eu\star$.  Multiplication of an endomorphism by $-1$ sends each
generalized $\lambda$-eigenspace to the generalized $-\lambda$-eigenspace,
without changing its dimension or Jordan-block sizes.  The remaining claims
follow immediately.
\end{proof}

The intrinsic Euler field satisfies
$\mathcal L_E(\circ)=\circ$ and $[e,E]=e$, and its canonical powers satisfy
the Witt identities used below; see \cite{DavidHertling}.  Lemma~\ref{lem:euler-sign}
allows us to apply those identities to the intrinsic $F$-manifold while
reading the maximal-eigenvalue locus from the A-model operator
$\Eu\star$ used by KKPY.  No coordinate formula for the Euler field away
from the small point is used.

\begin{lemma}[The Hodge-fixed base]\label{lem:hodge-base}
\coverage{fragment}%
\lean{hodgeFixedSubalgebra_closed_under_eulerMultiplication,%
  hodgeFixedSubalgebra_eigenvector_transfers,%
  hodgeFixedTangentSpace_span_and_finrank, truncatedHyperplaneAlgebra_model}%
\imports{KKPYY:maximal-f-bundle, KKPYY:sec-5.2.2}%
Let $B_X^{\Hod}$ be the Hodge-fixed locus of the maximal $A$-model
$F$-bundle of a smooth projective variety $X$.  In the notation of
\cite[Section~5.2.2]{KKPYY}, its tangent space is obtained from
\[
 H(X)^{\Hod}:=
 \bigoplus_i\bigl(H^{i,i}(X)\cap H^{2i}(X,\Q)\bigr)
\]
by extension to the Novikov coefficient field.  The restriction of the
$F$-manifold multiplication to $T B_X^{\Hod}$ is the corresponding
restriction of quantum multiplication.  In particular $B_X^{\Hod}$ is an
$F$-submanifold and its intrinsic Euler endomorphism is the restriction of
$\kappa$.

For a smooth cubic threefold, Lefschetz gives
\[
 H(X)^{\Hod}\otimes\kk
 =H^{\mathrm{even}}(X)\otimes\kk
 =\langle 1,P,P^2,P^3\rangle_{\kk}.
\]
Thus $B_X^{\Hod}$ is four-dimensional and, at the small hyperplane point,
its tangent Euler endomorphism is, up to the sign of
Lemma~\ref{lem:euler-sign}, Cai's four-dimensional even Euler-multiplication
matrix.
\end{lemma}

\begin{proof}
The maximal $A$-model $F$-bundle and its multiplication are
$\Hod$-equivariant in the construction of \cite{KKPYY}.  Therefore the
fixed tangent subbundle is closed under multiplication, contains the unit,
and is preserved by the Euler field; this is the induced $F$-manifold
structure on the smooth fixed locus.  For a cubic threefold,
$H^{2i}(X,\Q)=\Q P^i$ for $0\le i\le3$, so the displayed definition of
$H(X)^{\Hod}$ is exactly the even cohomology.  The final assertion follows
from Lemma~\ref{lem:euler-sign}.
\end{proof}


\subsection{The small-even cubic packet}\label{subsec:shared-cubic-packet}

The following calculation is shared by the atom argument of this section and
the framed-monodromy calculation of Section~\ref{sec:framed-monodromy}.  It
is entirely a small-connection calculation and is unconditional.

\begin{proposition}[Small-even cubic block reduction]\label{prop:cubic-block-data}
\coverage{fragment}%
\lean{cubicSmallEven_blockReduction,%
  cubicZeroBlock_modifiedResidue_indicialPolynomial}%
\imports{Cai:small-even-connection}%
Let \(X\subset\PP^4\) be a smooth cubic threefold, let \(P=H|_X\), choose
the line class \(\ell\) with \(P\cdot\ell=1\), and put \(q=Q^\ell\) and
\(r=(3q)^{1/2}\).  After adjoining \(r\) to the Novikov coefficient field,
the small even Euler multiplication has characteristic polynomial
\[
  T^2(T^2-108q),
\]
with eigenvalues \(6r,-6r,0,0\); the zero generalized eigenspace is a
single rank-two Jordan block.  Moreover a constant coefficient change of
basis followed by a unique normalized block-off-diagonal gauge
\(I+O(z)\), using only integral powers of the original loop coordinate,
puts the horizontal system into three blocks
\begin{equation}\label{eq:cubic-integral-block-reduction}
 \begin{aligned}
 z^2\partial_z\widetilde S_+&=(6r+O(z^2))\widetilde S_+,\\
 z^2\partial_z\widetilde S_-&=(-6r+O(z^2))\widetilde S_-,\\
 z^2\partial_z
 \begin{pmatrix}\widetilde S_3\\\widetilde S_4\end{pmatrix}
 &=\bigl[J_0+zD_0+z^2E_0+O(z^3)\bigr]
 \begin{pmatrix}\widetilde S_3\\\widetilde S_4\end{pmatrix},
 \end{aligned}
\end{equation}
where
\begin{equation}\label{eq:cubic-shared-block-data}
 J_0=\begin{pmatrix}0&2\\0&0\end{pmatrix},\qquad
 D_0=\begin{pmatrix}-19/18&0\\0&19/18\end{pmatrix},\qquad
 E_0=\begin{pmatrix}0&-14/(81r^2)\\-8/81&0\end{pmatrix}.
\end{equation}
\end{proposition}

\begin{proof}
Cai's small even horizontal system, in the ordered classical basis
\(1,P,P^2,P^3\), is
\begin{equation}\label{eq:cubic-small-even-connection}
 \begin{aligned}
 z^2\partial_zS&=(K_X+zG_X)S,\\
 K_X&=2\begin{pmatrix}
 0&6q&0&36q^2\\
 1&0&15q&0\\
 0&1&0&6q\\
 0&0&1&0
 \end{pmatrix},\\
 G_X&=\frac12\begin{pmatrix}
 3&0&0&0\\0&1&0&0\\0&0&-1&0\\0&0&0&-3
 \end{pmatrix}.
 \end{aligned}
\end{equation}
This is the displayed system in \cite[Section~3, p.~4]{Cai}.\footnote{In
arXiv:2608.01577v1, one isolated prose line states the two nonzero Euler
eigenvalues with a factor two smaller.  We follow Cai's displayed matrix,
its transformed block, and the subsequent scalar equations, which are
mutually consistent and give the values $\pm6r$ used here.}  The constant
change of basis
\begin{equation}\label{eq:cubic-constant-block-basis}
 C=\begin{pmatrix}
 6r^3&-6r^3&0&-7r^2\\
 7r^2&7r^2&-2r^2&0\\
 3r&-3r&0&1\\
 1&1&1&0
 \end{pmatrix},\qquad
 \det C=-486r^5,\qquad S=C\widehat S,
\end{equation}
puts the constant term into
\begin{equation}\label{eq:cubic-constant-block-form}
 C^{-1}K_XC=
 \begin{pmatrix}
 6r&0&0&0\\
 0&-6r&0&0\\
 0&0&0&2\\
 0&0&0&0
 \end{pmatrix}.
\end{equation}
In particular its characteristic and minimal polynomials are both
\(T^2(T^2-108q)\).

Partition the coordinates as \(1\mathbin{|}1\mathbin{|}2\).  There is a
unique formal gauge
\begin{equation}\label{eq:cubic-integral-block-gauge}
 A(z)=I+\sum_{n\ge1}A_nz^n,
 \qquad A_n\in\operatorname{Mat}_4(\C[r,r^{-1}]),
\end{equation}
whose positive coefficients are block-off-diagonal and for which
\(\widehat S=A(z)\widetilde S\) makes every coefficient block diagonal.  At
order \(n\) the required coefficient is the unique solution of a Sylvester
equation between blocks whose scalar spectra are distinct; the nilpotent
part of the zero block is only a nilpotent perturbation of the invertible
Sylvester operator.  Thus the gauge and its inverse use integral powers of
\(z\).  Direct multiplication gives exactly
\eqref{eq:cubic-integral-block-reduction} and
\eqref{eq:cubic-shared-block-data}.  This agrees with Cai's reduction
\cite[Section~3, pp.~4--5]{Cai}.
\end{proof}

\subsection{The cubic zero packet is an atom}

We first remove a point that cannot safely be delegated to a worked example:
the double zero eigenvalue at the small hyperplane point must remain a
repeated eigenvalue on the Hodge base, so that this point lies in the
maximal-eigenvalue locus used to define atoms.

\subsubsection{A discriminant lemma for a regular four-dimensional \texorpdfstring{$F$}{F}-manifold}

Let $(M,\circ,e,E)$ be an $F$-manifold over a characteristic-zero field.  Put
\[
  U=C_E:T_M\longrightarrow T_M,
  \qquad X_s=E^{\circ s},\quad X_0=e.
\]
At a point where $U$ is regular, $X_0,\dots,X_{n-1}$ form a local frame, where
$n=\dim M$.  The canonical frame satisfies
\begin{equation}\label{eq:witt}
  [X_i,X_j]=(j-i)X_{i+j-1}.
\end{equation}
These are the standard Witt identities for a regular $F$-manifold; see
\cite{DavidHertling}.  We need only the case $n=4$.

Write
\begin{equation}\label{eq:charpoly-lambda}
  P(T)=\det(TI-U)
  =T^4+\lambda_3T^3+\lambda_2T^2+\lambda_1T+\lambda_0
\end{equation}
and let $\Delta=\disc(P)$.

\begin{lemma}[Discriminant differential equation]\label{lem:disc}
\coverage{conditional_deduction}%
\lean{quarticDiscriminant_eq_squared_root_differences,%
  eulerFrame_discriminant_logarithmicDerivatives,%
  eulerFrame_discriminant_differential,%
  eulerFrame_discriminant_eq_zero_of_vanishing_at_base_point,%
  eulerFrame_data_nonempty}%
\imports{DavidHertling:euler-witt-identities}%
Assume $U$ is regular at $p\in M$.  On the germ of $M$ at $p$,
\begin{align}
  X_0(\Delta)&=0,\label{eq:disc0}\\
  X_1(\Delta)&=12\Delta,\label{eq:disc1}\\
  X_2(\Delta)&=-6\lambda_3\Delta,\label{eq:disc2}\\
  X_3(\Delta)&=(6\lambda_3^2-10\lambda_2)\Delta.\label{eq:disc3}
\end{align}
Consequently there is a regular one-form $\omega$ on the germ such that
\begin{equation}\label{eq:dDelta}
  d\Delta=\Delta\,\omega.
\end{equation}
If $\Delta(p)=0$, then $\Delta$ vanishes identically on the germ at $p$.
\end{lemma}

\begin{proof}
Cayley--Hamilton, applied to the unit, is the vector-field relation
\begin{equation}\label{eq:CH-vector}
 X_4+\lambda_3X_3+\lambda_2X_2+\lambda_1X_1+\lambda_0X_0=0.
\end{equation}
Taking Lie brackets of \eqref{eq:CH-vector} with $X_s$, using
\eqref{eq:witt}, reducing every $X_m$ with $m\ge4$ by
\eqref{eq:CH-vector} and its products by $E$, and comparing coefficients in
the frame $X_0,\dots,X_3$ gives the following coefficient identities:
\begin{align}
 X_0(\lambda_k)&=-(k+1)\lambda_{k+1},\label{eq:lambda0}\\
 X_1(\lambda_k)&=(4-k)\lambda_k,\label{eq:lambda1}\\
 X_2(\lambda_k)&=-\lambda_3\lambda_k+(5-k)\lambda_{k-1},\label{eq:lambda2}\\
 X_3(\lambda_k)&=(\lambda_3^2-2\lambda_2)\lambda_k
 +(6-k)\lambda_{k-2}-\lambda_3\lambda_{k-1},\label{eq:lambda3}
\end{align}
for $0\le k\le3$, with the conventions $\lambda_4=1$ and
$\lambda_j=0$ for $j<0$.  These are the $n=4$ cases of the coefficient
identities in \cite[(19),(20),(24),(25)]{DavidHertling}; the calculation above shows that
only the Witt identities and Cayley--Hamilton enter.

For completeness, we now compute the action on the discriminant without any
appeal to a splitting theorem on $M$.  Work universally in the polynomial
ring in formal roots $\mu_1,\dots,\mu_4$, so that
\[
  P(T)=\prod_{i=1}^4(T-\mu_i).
\]
The derivations of the coefficient ring specified by
\eqref{eq:lambda0}--\eqref{eq:lambda3} are the restrictions of
\[
  D_s=\sum_{i=1}^4\mu_i^s\frac{\partial}{\partial\mu_i},
  \qquad 0\le s\le3.
\]
This is checked directly on the elementary symmetric functions, and is
exactly equivalent to \eqref{eq:lambda0}--\eqref{eq:lambda3}.  Since
\[
  \Delta=\prod_{i<j}(\mu_i-\mu_j)^2,
\]
we obtain
\begin{equation}\label{eq:logdisc}
  \frac{D_s\Delta}{\Delta}
  =2\sum_{i<j}\frac{\mu_i^s-\mu_j^s}{\mu_i-\mu_j}.
\end{equation}
The right side is symmetric, hence is a polynomial in the $\lambda_k$.
For $s=0,1,2,3$ it equals respectively
\[
  0,\qquad 12,\qquad -6\lambda_3,
  \qquad 6\lambda_3^2-10\lambda_2.
\]
Indeed, the $s=2$ quotient is $\mu_i+\mu_j$, and the $s=3$ quotient is
$\mu_i^2+\mu_i\mu_j+\mu_j^2$; summing over pairs and using
$\sum_i\mu_i=-\lambda_3$ and
$\sum_{i<j}\mu_i\mu_j=\lambda_2$ gives the displayed expressions.  This
proves \eqref{eq:disc0}--\eqref{eq:disc3}.  Since
$X_0,\dots,X_3$ form a frame, these equations are equivalent to
\eqref{eq:dDelta} for a regular one-form $\omega$.

It remains to prove the vanishing statement in the non-archimedean analytic
germ.  Let $\mathfrak m_p$ be the maximal ideal of the local ring.  Passing
to the completed local ring at the rigid point $p$, choose regular parameters
$y_1,\dots,y_4$; smoothness gives a noncanonical isomorphism
\[
  \widehat{\OO}_{M,p}\simeq \kappa(p)[[y_1,\dots,y_4]],
\]
where $\kappa(p)$ is the residue field of $p$.  (In the cubic application
below, $p=b$ is $\kk$-rational.)
Suppose $\Delta\ne0$ there.  Since $\Delta(p)=0$, let $\Delta_m$ be its
lowest nonzero homogeneous part, of degree $m\ge1$.  In
$d\Delta=\Delta\omega$, the left side has lowest possible degree $m-1$,
whereas the right side has degree at least $m$.  Hence $d\Delta_m=0$.
Characteristic zero implies that all partial derivatives of $\Delta_m$
vanish, so Euler's identity gives
$m\Delta_m=\sum y_i\partial_{y_i}\Delta_m=0$, a contradiction.  Thus
$\Delta=0$ in the completion.  Since the local ring is Noetherian, Krull's
intersection theorem makes the map to its completion injective, so
$\Delta=0$ as an analytic germ.
\end{proof}

\subsubsection{Application to a cubic threefold}

Let $X\subset\PP^4$ be a smooth cubic threefold and use the notation of
Proposition~\ref{prop:cubic-block-data}.  That proposition gives the small
even characteristic polynomial
\begin{equation}\label{eq:cubic-charpoly}
  T^2(T^2-108q)
\end{equation}
and the eigenvalues
\begin{equation}\label{eq:cubic-eigs}
  6r,\quad -6r,\quad 0,\quad0,
\end{equation}
with a single rank-two Jordan block at zero.
By Lemma~\ref{lem:euler-sign}, the intrinsic $F$-manifold Euler
endomorphism is this operator up to an overall sign.  In either sign the minimal polynomial
equals the characteristic polynomial, so the Euler endomorphism is regular
at the small hyperplane point $b$, and the quartic discriminant has the same
zero locus.

By Lemma~\ref{lem:hodge-base}, the Hodge-fixed base
$B_X^{\Hod}$ is a four-dimensional $F$-submanifold and its tangent Euler
endomorphism is exactly the four-dimensional operator to which
Lemma~\ref{lem:disc} applies (up to the harmless overall sign of
Lemma~\ref{lem:euler-sign}).  Since the quartic
\eqref{eq:cubic-charpoly} has a repeated root, its discriminant vanishes at
$b$.  Hence it vanishes on an admissible neighborhood of $b$.

KKPY prove that $B_X^{\Hod}$ is connected and smooth.  It is therefore
connected and normal.  The zero locus of the discriminant is an analytic
subset containing a nonempty admissible open, so \cite[Lemma~2.1.4]{Conrad}
gives
\begin{equation}\label{eq:disc-global}
  \Delta\equiv0\qquad\text{on }B_X^{\Hod}.
\end{equation}
Thus the even Euler endomorphism has at most three distinct eigenvalues at
every point, while \eqref{eq:cubic-eigs} shows that it has exactly three at
$b$.

We must compare this count with the spectrum on all of $H^\bullet(X)$, since
KKPY define the maximal-eigenvalue locus using the full $F$-bundle.

\begin{lemma}[Spectrum transfer]\label{lem:spectrum-transfer}
\coverage{conditional_deduction}%
\lean{eulerMultiplication_eigenvalues_module_eq_algebra}%
At every even point of the Hodge base, the set of eigenvalues of Euler
multiplication on $H^\bullet(X)$ is the same as the set of eigenvalues on
$H^{\mathrm{even}}(X)$.
\end{lemma}

\begin{proof}
The even quantum cohomology $A=H^{\mathrm{even}}(X)$ is a finite-dimensional
commutative algebra over the coefficient field, and $H^\bullet(X)$ is a
unital $A$-module because the bulk point is even.  Let
$A=\bigoplus_i A_i$ be the generalized-eigenvalue decomposition for
multiplication by the Euler element, with corresponding orthogonal
idempotents $\varepsilon_i$.  Then
\[
  H^\bullet(X)=\bigoplus_i\varepsilon_iH^\bullet(X).
\]
Each summand is nonzero, since
$\varepsilon_i=\varepsilon_i\cdot1\in\varepsilon_iH^\bullet(X)$.  On the
local algebra $A_i$, the element $E-\lambda_i$ is nilpotent, so its action on
every finite $A_i$-module is nilpotent.  Hence the only eigenvalue on
$\varepsilon_iH^\bullet(X)$ is $\lambda_i$.  Every even eigenvalue therefore
occurs on the full module, and no additional eigenvalue occurs there.
\end{proof}

It follows from \eqref{eq:disc-global} and Lemma~\ref{lem:spectrum-transfer}
that the number of eigenvalues on the full cohomology is everywhere at most
three and is exactly three at $b$.  Therefore $b$ lies in the open locus
$U_X$ of KKPY where that number is globally maximal.

The generalized eigenbundle ranks are locally constant on the unramified
reduced spectral cover and hence constant on each connected component.  At
$b$ the two nonzero branches in \eqref{eq:cubic-eigs} have rank one.  The
zero branch has even rank two.  Its odd rank is ten: indeed
$b_3(X)=10$, since
\[
  c(TX)=\frac{(1+P)^5}{1+3P}=1+2P+4P^2-2P^3,
\]
so $\chi_{\mathrm{top}}(X)=\int_Xc_3(TX)=-6$, while Lefschetz and Poincar\'e
duality give
$b_0=b_2=b_4=b_6=1$ and $b_1=b_5=0$.  Thus $4-b_3=-6$.
Moreover Euler multiplication annihilates $H^3(X)$ at $b$.  In Novikov
degree $d\ge0$, multiplication by $c_1(X)$ sends a class of degree $3$ to
cohomological degree
\[
  3+2-2c_1(X)\cdot d=5-4d.
\]
For $d=0$ this is $H^5(X)=0$, for $d=1$ it is $H^1(X)=0$, and for $d\ge2$
it is negative.  Hence the entire $H^3(X)$ lies in the zero packet.

We have therefore proved the following without using KKPY's worked cubic
example.

\begin{proposition}[The cubic zero atom]\label{prop:cubic-atom}
\coverage{fragment}%
\lean{cubicThreefold_totalChernClass_identity, cubicThreefold_bettiThree_eq_ten,%
  cubicThreefold_eulerMultiplication_annihilates_degreeThree,%
  cubicZeroPacket_component_excludes_rankOne_branches,%
  cubicZeroPacket_parityRanks_eq_two_and_ten,%
  cubicThreefold_totalChernClass_uniqueness}%
\imports{KKPYY:def-5.21, KKPYY:prop-5.23}%
For every smooth cubic threefold $X$, the small hyperplane point
$b\in B_X^{\Hod}$ lies in the maximal-eigenvalue locus $U_X$.  Over $b$ the
unramified reduced spectral cover has two rank-one branches corresponding to
the two nonzero eigenvalues and one branch whose generalized eigenbundle has
parity ranks
\begin{equation}\label{eq:parity-ranks}
  (\rk_{\bar0},\rk_{\bar1})=(2,10).
\end{equation}
The rank-twelve branch is a connected component of the cover and hence
defines a Hodge atom, denoted $\alpha_X$.
\end{proposition}

\begin{proof}
Only the last assertion remains.  Over $U_X$ the reduced spectral cover is
finite etale.  On each connected component of that cover, the corresponding
generalized-eigenbundle is locally free of constant rank.  The fibre over
$b$ consists of three spectral points: the point belonging to the zero packet
has generalized-eigenbundle rank $12$, while the two points belonging to the
nonzero eigenvalues have rank $1$.  Hence neither rank-one point can lie on
the connected component through the rank-twelve point.  The fibre of that
component over $b$ therefore consists only of the zero spectral point; in
particular the component is the atomic component determined by the zero
packet (and has degree one over its image in a neighborhood of $b$).  Its
parity ranks are \eqref{eq:parity-ranks}.
\end{proof}

By KKPY's invariance of the $\Hod$-representation of an atom, the parity ranks
\eqref{eq:parity-ranks} are invariants of the global Hodge atom $\alpha_X$.


\subsection{A rank-two atomic residue discriminant}

We now construct the invariant that distinguishes $\alpha_X$ from every
low-dimensional atom that could have the same parity ranks.

\subsubsection{Pairing and spectral factors}

Let $(\mathcal H,\nabla)/B$ be the non-archimedean $A$-model $F$-bundle of a
smooth projective variety, restricted to the purely even Hodge-fixed base.
In the non-archimedean construction of \cite[Section~3.5.2]{KKPYY}, the
inverse Gram matrix of the Poincar\'e form is a constant analytic section
$\mathrm{PD}^{-1}$ and is used to define quantum multiplication.  Thus the
Poincar\'e form itself is available on this analytic $F$-bundle, constant in
the standard cohomology frame.  We prove its horizontality below directly
and coefficientwise, rather than importing the later convergent
complex-analytic discussion of the pairing.  Define the $u$-sesquilinear pairing
\[
  \psi_u(s,t)=(s(u),t(-u))_{\mathrm{Poinc}}.
\]

\begin{lemma}[Algebraic horizontality]\label{lem:horiz}
\coverage{fragment}%
\lean{quantumPairing_horizontality_of_selfAdjoint_multiplication}%
\imports{KKPYY:sec-3.5.2}%
The pairing $\psi$ is horizontal for the non-archimedean quantum connection.
No convergence assumption is required.
\end{lemma}

\begin{proof}
Quantum multiplication is Frobenius:
\[
  (a\star b,c)=(a,b\star c).
\]
Thus every quantum multiplication operator, in particular Euler
multiplication and multiplication by an even tangent vector, is self-adjoint
for the Poincar\'e pairing.  The grading operator
\[
  \Gr|_{H^a}=\frac{a-\dim X}{2}\,\id
\]
is anti-self-adjoint because Poincar\'e duality pairs $H^a$ with
$H^{2\dim X-a}$.  Substituting these two identities into the displayed
quantum-connection formulas proves horizontality coefficientwise in the
Novikov and bulk variables.  The proof is algebraic and therefore takes
place directly over the non-archimedean coefficient ring.
\end{proof}

The even part is a sub-$F$-bundle: the Hodge-fixed base is purely even and
quantum multiplication by an even class, the Euler operator, and the grading
operator all preserve cohomological parity.

We also need that the pairing restricts nondegenerately to a spectral factor,
without assuming an isometric block-splitting gauge.

\begin{lemma}[Orthogonality of separated spectral factors]\label{lem:orthogonal}
\coverage{conditional_deduction}%
\lean{separatedSpectralFactors_pairing_blockDiagonal,%
  separatedSpectralFactors_pairing_eq_zero_of_horizontality,%
  separatedSpectralFactors_labelledPairing_blockDiagonal_of_horizontality,%
  separatedSpectralFactors_evenPart_pairing_nondegenerate,%
  separatedSpectralFactors_equalEigenvalues_admit_nonzero_horizontalPairing}%
Suppose locally the $F$-bundle is split into factors whose leading Euler
eigenvalues $\lambda_i$ are pairwise distinct.  Then the horizontal Poincar\'e
pairing is block diagonal for that splitting; in particular its restriction
to each factor, and to the even part of each factor, is nondegenerate.
\end{lemma}

\begin{proof}
Choose any regular block-diagonalizing frame for the connection and write the
pairing matrix as $P(u)=\sum_{m\ge0}P_m u^m$.  On the $(i,j)$ off-diagonal
block, the coefficient of $u^{-1}$ in horizontality is
\[
  U_i^T(P_0)_{ij}-(P_0)_{ij}U_j=0.
\]
The Sylvester operator
\[
  X\longmapsto U_i^TX-XU_j
\]
is invertible because $\lambda_i-\lambda_j$ is a unit and the remaining parts of
$U_i,U_j$ are nilpotent.  Hence $(P_0)_{ij}=0$.  Inductively, the coefficient
of each successive power of $u$ has the same invertible Sylvester operator
on $(P_m)_{ij}$ and a right-hand side involving only earlier off-diagonal
coefficients, already zero.  Thus $P_{ij}(u)=0$ identically.  Since the full
pairing is nondegenerate, every diagonal restriction is nondegenerate.  The
Poincar\'e form pairs only equal parity, so the even restriction is also
nondegenerate.
\end{proof}

\subsubsection{The elementary modification}

Let $\mathcal U$ be a connected component of the unramified spectral cover
representing a geometric Hodge atom, and suppose the even part $E$ of its
geometric atomic $F$-bundle has rank two.  Let $\mathsf U$ be its leading
Euler operator.  On the spectral component $\mathsf U$ has a single
eigenvalue $\lambda$; put
\begin{equation}\label{eq:Ndef}
  N=\mathsf U-\lambda I,\qquad \lambda=\frac12\tr(\mathsf U).
\end{equation}
Then $N^2=0$.  Assume $N$ is generically nonzero on $\mathcal U$.

Center the scalar irregular part algebraically.  If the original matrices
are $A^{\rm old}$ and $B^{\rm old}_\delta$, replace them by
\[
  A=A^{\rm old}-u^{-1}\lambda I,
  \qquad
  B_\delta=B^{\rm old}_\delta+u^{-1}(\delta\lambda)I.
\]
This is the connection obtained formally by multiplying horizontal sections by
$\exp(\lambda/u)$, but no exponential is adjoined and no convergence is
being asserted.  The added scalar one-form is exact, so flatness is
unchanged.  In a local frame horizontal sections for the centered connection
satisfy
\begin{equation}\label{eq:localpair}
  u\partial_u y=A(u)y,\qquad
  \delta y=B_\delta(u)y,
\end{equation}
with
\begin{align}
  A(u)&=u^{-1}N+A_0+uA_1+u^2A_2+\cdots,\label{eq:Aexp}\\
  B_\delta(u)&=u^{-1}C_\delta+C_{\delta,0}
       +uC_{\delta,1}+\cdots.\label{eq:Bexp}
\end{align}
The centering does not alter the pairing identity: the scalar terms in the
$u$ and $-u$ factors have opposite signs and cancel.  Flatness is
\begin{equation}\label{eq:flatAB}
  \delta A-u\partial_uB_\delta+[A,B_\delta]=0.
\end{equation}

On the open locus $\mathcal U^\times=\{N\ne0\}$,
\[
  L:=\im N=\ker N
\]
is a line.

\begin{lemma}[The regular coefficient preserves the nilpotent line]
\label{lem:A0preserve}
\coverage{conditional_deduction}%
\lean{atomicRankTwo_regularCoefficient_preserves_nilpotentLine,%
  atomicRankTwo_pairingEquations_of_loopHorizontality}%
On $\mathcal U^\times$,
\begin{equation}\label{eq:A0L}
  A_0L\subset L.
\end{equation}
\end{lemma}

\begin{proof}
By Lemmas~\ref{lem:horiz} and \ref{lem:orthogonal}, the rank-two even factor
has a horizontal nondegenerate pairing
\[
  P(u)=P_0+uP_1+u^2P_2+\cdots,
\]
with $P_0$ symmetric.  Horizontality in the $u$-direction reads
\begin{equation}\label{eq:pair-horizontal}
  u\partial_uP(u)+A(u)^TP(u)+P(u)A(-u)=0.
\end{equation}
The $u^{-1}$ coefficient gives
\begin{equation}\label{eq:Nself}
  N^TP_0=P_0N.
\end{equation}
Hence $N$ is self-adjoint for $P_0$.  Since $N^2=0$,
\[
  (Nx,Ny)=(x,N^2y)=0,
\]
so $L=\im N$ is isotropic.  In a nondegenerate two-dimensional symmetric
space, $L^\perp=L$.

The constant coefficient of \eqref{eq:pair-horizontal} is
\begin{equation}\label{eq:constant-pair}
 A_0^TP_0+P_0A_0+N^TP_1-P_1N=0.
\end{equation}
For $0\ne x\in L=\ker N$, sandwiching
\eqref{eq:constant-pair} between $x$ on both sides kills the two terms
containing $N$ and gives $2(A_0x,x)=0$.  Thus
$A_0x\in x^\perp=L^\perp=L$, proving \eqref{eq:A0L}.
\end{proof}

Define the elementary modification
\begin{equation}\label{eq:Emod}
  E^\sharp=\{s\in E:s\bmod u\in L\}.
\end{equation}
This is intrinsic.  Choose locally a frame $e_1,e_2$ with
\begin{equation}\label{eq:adapted}
  L=\langle e_1\rangle,\qquad
  Ne_1=0,
  \qquad Ne_2=\nu e_1,
  \qquad \nu\in\OO^\times.
\end{equation}
Then $E^\sharp$ has frame $e_1,ue_2$, corresponding to
$S=\operatorname{diag}(1,u)$.  The transformed $u$-matrix is
\begin{equation}\label{eq:Asharp}
  A^\sharp=S^{-1}AS-S^{-1}u\partial_uS.
\end{equation}
Because
\[
  S^{-1}E_{12}S=uE_{12},\qquad
  S^{-1}E_{21}S=u^{-1}E_{21},
\]
the term $u^{-1}N=\nu u^{-1}E_{12}$ becomes the regular term
$\nu E_{12}$.  The only possible pole coming from $A_0$ is its
$E_{21}$ coefficient, and Lemma~\ref{lem:A0preserve} says that coefficient
vanishes.  Every $u^mA_m$, $m\ge1$, remains regular after losing at most one
power of $u$.  Hence
\begin{equation}\label{eq:Rexp}
  A^\sharp(u)=R+uR_1+u^2R_2+\cdots.
\end{equation}
The conjugacy class of the residue $R$ is intrinsic to the modified lattice.

We use only its discriminant
\begin{equation}\label{eq:res-disc}
  \delta^\sharp
  :=(\tr R)^2-4\det R.
\end{equation}
This is invariant under conjugacy and also under adding a scalar multiple of
the identity to $R$.  The latter observation makes the invariant insensitive
to harmless scalar residue normalizations.

\subsubsection{Flatness freezes the residue discriminant}

\begin{proposition}[Rank-two rigidity]\label{prop:rank2-rigidity}
\coverage{fragment}%
\lean{atomicRankTwo_modifiedBase_pole_eq_zero,%
  atomicRankTwo_residueDiscriminant_constant_along_base}%
On $\mathcal U^\times$, the modified connection is regular in every base
direction and the residue satisfies
\begin{equation}\label{eq:Lax}
  \delta R=[G_\delta,R]
\end{equation}
for a regular matrix $G_\delta$.  Consequently
\begin{equation}\label{eq:dresdisc}
  \delta(\delta^\sharp)=0
\end{equation}
for every base vector field $\delta$.
\end{proposition}

\begin{proof}
The $u^{-2}$ coefficient of \eqref{eq:flatAB} is
\begin{equation}\label{eq:commCN}
  [N,C_\delta]=0.
\end{equation}
The $u^{-1}$ coefficient is
\begin{equation}\label{eq:u-1flat}
  \delta N+C_\delta+[N,C_{\delta,0}]+[A_0,C_\delta]=0.
\end{equation}
Since $\tr N=0$, taking traces gives $\tr C_\delta=0$.  On
$\mathcal U^\times$ the commutant of a nonzero rank-one nilpotent
$2\times2$ matrix is
$\OO I\oplus\OO N$.  Equations \eqref{eq:commCN} and
$\tr C_\delta=0$ therefore give
\begin{equation}\label{eq:CqN}
  C_\delta=q_\delta N.
\end{equation}

Transform the base equation to the modified lattice.  In the moving adapted
frame \eqref{eq:adapted}, the frame-connection contribution is included in
$C_{\delta,0}$; the matrix $S=\diag(1,u)$ itself is base-independent.  The
leading term $u^{-1}q_\delta N$ becomes regular.  Thus the only possible pole
has the form
\begin{equation}\label{eq:Bsharp}
  B_\delta^\sharp=u^{-1}K_\delta+G_\delta+O(u),
  \qquad K_\delta=k_\delta E_{21}.
\end{equation}
Flatness in the modified lattice gives at order $u^{-1}$
\begin{equation}\label{eq:Kflat}
  K_\delta+[R,K_\delta]=0.
\end{equation}
The $(1,2)$ entry of $R$ is exactly the unit $\nu$ coming from $N$, so write
\[
  R=\begin{pmatrix}a&\nu\\ c&d\end{pmatrix}.
\]
Then
\[
  [R,k_\delta E_{21}]
  =k_\delta
    \begin{pmatrix}\nu&0\\ d-a&-\nu\end{pmatrix}.
\]
The two diagonal entries of \eqref{eq:Kflat} are therefore
$\nu k_\delta$ and $-\nu k_\delta$.  Since $\nu$ is a unit,
$K_\delta=0$.  Hence the base connection is regular after the modification.
The constant coefficient of flatness is now
\[
  \delta R+[R,G_\delta]=0,
\]
which is \eqref{eq:Lax}.  Trace and determinant are invariant under
infinitesimal conjugation, so \eqref{eq:dresdisc} follows.  In particular
$\delta^\sharp$ is locally constant on $\mathcal U^\times$: on a smooth
rigid chart with regular parameters $t_1,\ldots,t_m$, apply
\eqref{eq:dresdisc} to the coordinate derivations.  In the completed local
ring every positive-degree coefficient of the resulting power series is
then zero, so the function is constant on each connected chart.
\end{proof}

\begin{proposition}[Exponent interpretation of the residue discriminant]
\label{prop:residue-discriminant-exponents}
\coverage{fragment}%
\lean{residueDiscriminant_eq_squared_eigenvalue_separation}%
\imports{PutSinger:ch-3}%
Fix a point of $\mathcal U^\times$.  Let $K$ be its coefficient field and
work in an algebraic closure $\overline K$ over which the residue $R$ splits.
Let $\rho_1,\rho_2$ be the eigenvalues of the residue of the canonical
modified lattice $E^\sharp$, counted with algebraic multiplicity.  Then:
\begin{enumerate}[label=(\roman*)]
\item the classes $[\rho_1],[\rho_2]\in\overline K/\Z$ are the formal
  exponent classes of the original \emph{centered} rank-two differential
  module; and
\item
\begin{equation}\label{eq:res-disc-exponents}
  \delta^\sharp=(\rho_1-\rho_2)^2.
\end{equation}
\end{enumerate}
Thus $\delta^\sharp$ records the squared separation of the residue
representatives selected by the canonical elementary modification.  It is
invariant under permuting the two representatives and under a common scalar
shift.  It is not, in general, determined by the two exponent classes modulo
$\Z$ alone.
\end{proposition}

\begin{proof}
The elementary modification is intrinsic, and in an adapted frame it is
implemented by $S=\diag(1,u)$ (or by $\diag(u,1)$ after reversing the
orientation of the nilpotent Jordan block).  In either form
$S\in\operatorname{GL}_2(\overline K((u)))$ and uses only integral powers of
the original loop coordinate.  Hence it is a meromorphic gauge isomorphism
of differential modules; because this particular gauge multiplies the
adapted basis vectors only by integral powers of $u$, it changes the
corresponding exponent representatives only by integers and therefore
preserves their classes in $\overline K/\Z$.

By Proposition~\ref{prop:rank2-rigidity}, the modified connection is regular
singular and has the form
\[
  u\partial_u y^\sharp=(R+O(u))y^\sharp.
\]
For a regular-singular formal differential module, the formal exponent
classes are the residue eigenvalues modulo $\Z$, including in the resonant
or nonsemisimple case; see, for example, the regular-singular part of the
formal classification in \cite[Chapter~3]{PutSinger}.  Therefore
$[\rho_1],[\rho_2]$ are the formal exponent classes of the original centered
block.

Finally,
\[
 (\tr R)^2-4\det R
 = (\rho_1+\rho_2)^2-4\rho_1\rho_2
 = (\rho_1-\rho_2)^2,
\]
which proves \eqref{eq:res-disc-exponents}.
\end{proof}

\subsubsection{Crossing the locus where \texorpdfstring{$N$}{N} vanishes}

The preceding calculation was made only where $N\ne0$.  That is enough.  We
first record the gluing point implicit in the spectral decomposition.

\begin{lemma}[Gluing of a separated spectral factor]\label{lem:factor-glue}
\coverage{fragment}%
\lean{atomicFactor_residueDiscriminant_invariant_under_blockDiagonal_frame_change}%
\imports{KKPYY:prop-5.23}%
Over the unramified locus $U_X$, the local $F$-bundle factors belonging to a
fixed connected component $\mathcal U$ of the reduced spectral cover glue,
up to regular block-diagonal change of frame.  In particular any
regular-gauge-invariant meromorphic expression formed from the factor is a
well-defined meromorphic function on $\mathcal U$.
\end{lemma}

\begin{proof}
At $u=0$, the proof of \cite[Proposition~5.23]{KKPYY} constructs the
generalized-eigenspace subbundle over $U_X$ associated with each component
of the reduced spectral cover, and the spectral decomposition theorem of
\cite{KKPYY} supplies the corresponding local analytic $F$-bundle factors.
Thus the factor itself is part of the ordinary atom package.  We record only
the uniqueness needed to glue the higher $u$-jet expressions used below.

Fix the splitting at $u=0$.  Among regular gauges $g=I+O(u)$ carrying the
connection to block-diagonal form, impose the normalization that every
positive $u$-coefficient of $g$ is block-off-diagonal.  At each order the
unknown off-diagonal coefficient is determined by a Sylvester operator of
the form
\[
 X\longmapsto U_iX-XU_j.
\]
It is invertible because the scalar parts of $U_i$ and $U_j$ are distinct on
$U_X$; the remaining nilpotent terms give only a nilpotent perturbation of
multiplication by their unit difference.  Hence a normalized splitting, if
chosen, is unique.  The analytic splittings supplied by KKPY therefore agree
on overlaps after a regular block-diagonal change of frame.  Consequently
regular-gauge-invariant meromorphic expressions formed from one spectral
factor glue along its connected component $\mathcal U$.
\end{proof}

\begin{lemma}[Meromorphic continuation on a spectral component]
\label{lem:crossing}
\coverage{fragment}%
\lean{atomicResidueDiscriminant_constant_on_formal_germ}%
\imports{Conrad:lem-2.1.4}%
Assume $N$ is generically nonzero on the connected spectral component
$\mathcal U$.  Then $\delta^\sharp$ extends as a meromorphic function on
$\mathcal U$ and is constant there.  In particular it is independent of the
chosen base point representing the same local atom.
\end{lemma}

\begin{proof}
The reduced spectral cover is unramified over the smooth Hodge base on
$U_X$.  Hence $\mathcal U$ is smooth and therefore normal.  A connected
normal rigid space is irreducible by the rigid-analytic irreducibility theory
of \cite{Conrad}.

By Lemma~\ref{lem:factor-glue}, the rank-two factor is defined along
$\mathcal U$ up to regular gauge.  Over its meromorphic function field, the
nonzero nilpotent matrix $N$ has rank one and its image line is generated by
any nonzero column.  Choose a meromorphic adapted frame.  In that frame the
coefficients of the connection and therefore the residue $R$ of the modified
lattice are meromorphic.  Since \eqref{eq:res-disc} is gauge invariant,
$\delta^\sharp$ is a meromorphic function on $\mathcal U$.

Choose a nonempty admissible open in $\mathcal U^\times$.  By
Proposition~\ref{prop:rank2-rigidity}, $\delta^\sharp$ is constant on a
smaller connected open there; call the value $c$.  The difference
$\delta^\sharp-c$ is an element of the meromorphic function field of the
integral rigid space $\mathcal U$.  To spell out the identity principle,
choose a connected normal affinoid subdomain meeting that open and write the
difference there as $a/b$ with $a,b$ analytic and $b$ nonzero.  On a
smaller nonempty admissible open where $b$ is invertible the quotient
vanishes, so $a$ vanishes there.  Conrad's identity principle for connected
normal rigid spaces gives $a=0$ on the affinoid.  Since a meromorphic
function on the irreducible space $\mathcal U$ is determined by its value
on any nonempty admissible open, $\delta^\sharp-c=0$ in its meromorphic
function field.  In particular the same value is obtained on every
component of $\mathcal U^\times$, even if those components are separated
by points at which $N=0$.
\end{proof}

\subsubsection{Descent to ordinary Hodge atoms}

\begin{proposition}[Atomic residue discriminant]\label{prop:atom-invariant}
\coverage{fragment}%
\lean{atomicResidueDiscriminant_invariant_under_frame_change}%
\imports{KKPYY:prop-5.22}%
Let $\alpha$ be an ordinary Hodge atom whose associated geometric atomic
$F$-bundle has even rank two and whose centered leading operator is
generically nonzero nilpotent.  Then the number $\delta^\sharp(\alpha)$
defined by \eqref{eq:res-disc} is a well-defined invariant of the ordinary
Hodge atom.
\end{proposition}

\begin{proof}
First consider an isomorphism of geometric atomic $F$-bundles.  It preserves
cohomological parity, hence preserves the even rank.  Its value at $u=0$
conjugates the centered leading operator $N$ to $N'$; consequently generic
nonvanishing of $N$ is preserved, and the isomorphism carries
$L=\im N$ to $L'=\im N'$.  Hence it carries the intrinsic elementary
modification $E^\sharp$ to $(E')^\sharp$.  The residues are conjugate after
applying the same trace-centering on both sides.  Any residual scalar shift
of $R$ is in any event killed by \eqref{eq:res-disc}; therefore both the
domain condition and $\delta^\sharp$ are preserved.

Second, changing the base point within one connected component of the
unramified spectral cover does not change the even rank or the property that
$N$ is generically nonzero on that component, and it preserves
$\delta^\sharp$ by Lemma~\ref{lem:crossing}.  Thus the domain of definition
and the value of $\delta^\sharp$ are invariant under the two equivalences
used by KKPY for geometric atomic $F$-bundles.

There are two further quotients in the passage to global Hodge atoms.  First,
KKPY quotient components attached to a fixed variety by $\operatorname{Aut}(X)$.
An automorphism pulls back the $A$-model $F$-bundle and carries the relevant
spectral component to the corresponding one, so the isomorphism case above
shows that both the domain condition and $\delta^\sharp$ are unchanged.
Second, $\HAtoms$ is further quotiented by the elementary relations coming
from disjoint unions, blowups with smooth centers, and projective bundles.
Proposition~5.22 of \cite{KKPYY} proves exactly that the map from these local
atom classes to geometric atomic $F$-bundle equivalence classes factors
through those elementary relations: the two sides of each relation have the
same image in the target set of geometric atomic $F$-bundle classes.  Since
the rank-two/generic-nilpotence domain condition and $\delta^\sharp$ have
already been shown to be well defined on that target set, composing with the
map of Proposition~5.22 gives a well-defined invariant on precisely the
ordinary Hodge-atom classes satisfying the stated domain condition.
\end{proof}

\begin{remark}
The Poincar\'e pairing is used only as an auxiliary structure to prove the
intrinsic property \eqref{eq:A0L} of the unenhanced atomic connection.  No
pairing is added to the definition of the atom, and morphisms of atoms are
not required to preserve an additional enhancement.
\end{remark}


\subsection{The cubic residue and low-dimensional exclusion}

\subsubsection{The residue discriminant of the cubic atom}
For the specialization of the abstract rank-two construction above to
Proposition~\ref{prop:cubic-block-data}, we identify its formal loop
coordinate $u$ with the original loop coordinate $z$ used in the cubic
small-even system.  By Proposition~\ref{prop:cubic-block-data}, the
rank-two zero block is
\[
 z^2\partial_zY=(J_0+zD_0+z^2E_0+O(z^3))Y
\]
with the matrices in \eqref{eq:cubic-shared-block-data}.  Dividing by \(z\),
the centered leading term is \(N=J_0=2E_{12}\) and
\(L=\langle e_1\rangle\).  The canonical elementary modification uses
\(S=\diag(1,z)\), so Formula~\eqref{eq:Asharp} gives the residue
\begin{equation}\label{eq:RX}
 R_X
 =2E_{12}+D_0-\frac8{81}E_{21}-\diag(0,1)
 =\begin{pmatrix}
  -19/18&2\\
  -8/81&1/18
 \end{pmatrix}.
\end{equation}
Hence
\[
  \tr R_X=-1,\qquad \det R_X=\frac5{36},
\]
and its residue characteristic polynomial is
\[
 T^2+T+\frac5{36}
 =\left(T+\frac16\right)\left(T+\frac56\right).
\]
Thus the canonical modified-residue representatives are
$\rho_1=-1/6$ and $\rho_2=-5/6$.  Proposition~\ref{prop:residue-discriminant-exponents}
therefore gives
\begin{equation}\label{eq:cubic-delta}
 \boxed{\delta^\sharp(\alpha_X)
   =\left(-\frac16+\frac56\right)^2=\frac49.}
\end{equation}
The occurrence of \((E_0)_{21}=-8/81\) in \(R_X\) is important: the
modification retains precisely the next coefficient that controls this
formal exponent separation.

\subsubsection{No curve represents the cubic atom}

Suppose $\alpha_X$ is represented by a smooth projective curve $C$.
If $C=\PP^1$, KKPY's projective-bundle formula gives
\[
  \CF(\PP^1)=2\CF(\mathrm{pt}),
\]
so only point atoms occur; their rank is one, whereas $\alpha_X$ has rank
$12$.

Assume $g(C)\ge1$.  Then $K_C$ is nef, so KKPY's nef-canonical lemma says
that $C$ has a single Hodge atom, carrying all of $H^\bullet(C)$.  Its parity
ranks are $(2,2g)$.  By the Hodge-representation invariance of atoms and
\eqref{eq:parity-ranks}, an identification with $\alpha_X$ would force
\begin{equation}\label{eq:g5}
  g=5.
\end{equation}
We now exclude that possibility by the residue discriminant.

Let $p\in H^2(C)$ be normalized by $\int_Cp=1$ and put $a=2-2g$.
For $g\ge1$ there are no nonconstant genus-zero maps $\PP^1\to C$, so the
big genus-zero quantum product is the classical cup product.  On the even
basis $1,p$, after removing the scalar $H^0$-direction, Cai's connection is
\begin{equation}\label{eq:curve-connection}
 u\partial_uY=
 \left[
  u^{-1}aE_{21}+\begin{pmatrix}1/2&0\\0&-1/2\end{pmatrix}
 \right]Y;
\end{equation}
see \cite[\S4]{Cai}.  For $g>1$, $a\ne0$ and
$N=aE_{21}$, so $L=\langle e_2\rangle$.  The elementary modification now
uses $S=\diag(u,1)$ and gives
\begin{equation}\label{eq:RC}
  R_C
  =aE_{21}+\diag(1/2,-1/2)-\diag(1,0)
  =aE_{21}-\frac12I.
\end{equation}
The modified residue has the repeated eigenvalue $-1/2$, so its two
canonical residue representatives coincide.  By
Proposition~\ref{prop:residue-discriminant-exponents},
\begin{equation}\label{eq:curve-delta}
  \boxed{\delta^\sharp(C)
    =\left(-\frac12+\frac12\right)^2=0.}
\end{equation}
Equivalently, its modified residue characteristic polynomial is
$(T+1/2)^2$.

For the only possible competitor $g=5$, equations
\eqref{eq:cubic-delta} and \eqref{eq:curve-delta} contradict the
well-definedness of $\delta^\sharp$ on ordinary Hodge atoms.  We have proved:

\begin{proposition}\label{prop:no-curve}
\coverage{conditional_deduction}%
\lean{cubicAtom_not_represented_by_curve}%
\imports{Cai:curve-calculation}%
The cubic atom $\alpha_X$ is not represented by any smooth projective
variety of dimension at most one.
\end{proposition}

\subsubsection{No surface represents the cubic atom}

\begin{proposition}\label{prop:no-surface}
\coverage{conditional_deduction}%
\lean{cubicAtom_not_represented_by_surface}%
\imports{KKPYY:lem-5.24, KKPYY:prop-5.30}%
The atom $\alpha_X$ is not represented by a smooth projective surface.
Together with Proposition~\ref{prop:no-curve} and additivity for disjoint
unions, this excludes every smooth projective variety of dimension at most
two.  Consequently
\begin{equation}\label{eq:not-lowdim}
  \alpha_X\notin\HAtoms_{\dim\le2}.
\end{equation}
\end{proposition}

\begin{proof}
Suppose $\alpha_X$ occurs in the atomic composition of a smooth projective
surface $S$.  Blow down $(-1)$-curves until reaching a minimal model
$S_{\min}$.  KKPY's blowup chemical formula for a point blowup is
\[
  \CF(\Bl_pS)=\CF(S)+\CF(p).
\]
A point atom has rank one, whereas $\alpha_X$ has rank twelve.  Therefore
$\alpha_X$ already occurs on $S_{\min}$.

If $K_{S_{\min}}$ is nef, KKPY's nef-canonical lemma says that
$S_{\min}$ has a single atom whose underlying $\Hod$-representation is all
of $H^\bullet(S_{\min})$.  Its even part has dimension at least three:
$H^0$, an ample class in $H^2$, and $H^4$ are linearly independent.  This
contradicts the even rank $2$ of $\alpha_X$.

If $K_{S_{\min}}$ is not nef, the classification of minimal smooth complex
projective surfaces gives either $\PP^2$ or a geometrically ruled surface
$\PP_C(V)$.  By KKPY's projective-bundle formula,
\[
  \CF(\PP^2)=3\CF(\mathrm{pt}),
  \qquad
  \CF(\PP_C(V))=2\CF(C).
\]
Thus every atom occurring there is a point or curve atom.  Proposition
\ref{prop:no-curve} excludes $\alpha_X$ in both cases.  This proves
\eqref{eq:not-lowdim}.
\end{proof}


\subsection{Proof of the one-stabilization theorem}

We can now finish without transporting framed monodromy through a birational
factorization.

\begin{proof}[Proof of Theorem~\ref{thm:every-cubic}]
Since
\[
  X\times\PP^1=\PP_X(\OO_X^{\oplus2}),
\]
KKPY's ordinary projective-bundle chemical formula gives
\begin{equation}\label{eq:product-CF}
  \CF(X\times\PP^1)=2\CF(X).
\end{equation}
Hence the Hodge atom \(\alpha_X\) occurs with positive multiplicity in the
atomic composition of the smooth projective fourfold \(X\times\PP^1\).

KKPY's ordinary Hodge-atom non-rationality criterion
\cite[Proposition~5.30]{KKPYY} states that if a smooth projective \(d\)-fold
contains an atom not realizable by any smooth projective variety of
dimension at most \(d-2\), then the \(d\)-fold is not rational.  Here
\(d=4\), and Proposition~\ref{prop:no-surface} gives
\[
  \alpha_X\notin\HAtoms_{\dim\le2}.
\]
Therefore \(X\times\PP^1\) is irrational.
\end{proof}

\begin{remark}[Scope of the atom obstruction]
The argument is specific to one stabilization.  For \(X\times\PP^m\) with
\(m\ge2\), the ordinary atom criterion permits representatives in dimension
at least three, and the cubic atom already has the threefold \(X\) as such a
representative.  Stronger stabilization statements therefore require a
finer decoration or a different invariant.
\end{remark}
